\chapter{Transmission Lines, Reflection, and SWR}
\label{ch:lines}
\section{Why Lines Matter}
Transmission lines are distributed systems. When a feed line is a significant
fraction of a wavelength, voltage and current vary along the line and a lumped
wire approximation fails.

\begin{controlsbox}
A transmission line is a distributed-parameter dynamical system rather than a
lumped ODE. A lumped RLC circuit has a finite state vector; a line has a
continuum of states \(v(x,t)\) and \(i(x,t)\) indexed by position \(x\). Position
becomes part of the state description, so the model is a partial differential
equation with an infinite set of natural modes rather than a handful of
eigenvalues.
\end{controlsbox}

\section{The Line as a Distributed System: Telegrapher's Equations}
Model an infinitesimal slice \(\dd x\) of a lossless line as a series inductance
\(L'\,\dd x\) and a shunt capacitance \(C'\,\dd x\), where \(L'\) and \(C'\) are
inductance and capacitance per unit length. Applying KVL and KCL to the slice
gives the \emph{telegrapher's equations},
\[
\frac{\partial v}{\partial x}=-L'\frac{\partial i}{\partial t},
\qquad
\frac{\partial i}{\partial x}=-C'\frac{\partial v}{\partial t}.
\]
Differentiating one and substituting the other eliminates \(i\) and yields the
\emph{wave equation}:
\[
\frac{\partial^2 v}{\partial x^2}=L'C'\,\frac{\partial^2 v}{\partial t^2}.
\]
Its solutions are waves traveling in \(\pm x\) at phase velocity
\[
v_p=\frac{1}{\sqrt{L'C'}},
\]
and the ratio of voltage to current in a single traveling wave is the
characteristic impedance
\[
Z_0=\sqrt{\frac{L'}{C'}}.
\]
So \(Z_0\) and the velocity factor are not empirical constants---they fall out of
the per-unit-length geometry of the line.

\section{Wavelength and Velocity Factor}
\[
\lambda=\frac{v_p}{f}=\frac{VF\,c}{f}.
\]
Velocity factor \(VF=v_p/c\) is the ratio of wave speed in the line to the speed
of light. Because \(v_p=1/\sqrt{L'C'}\) and a dielectric raises \(C'\), a line
with dielectric has \(VF<1\); its physical length is shorter than the free-space
electrical length.

\section{Characteristic Impedance}
\(Z_0=\sqrt{L'/C'}\) is the voltage-to-current ratio of a traveling wave. Common
coaxial systems use \(\SI{50}{\ohm}\); \(\SI{75}{\ohm}\), \(\SI{300}{\ohm}\), and
\(\SI{450}{\ohm}\) appear in specific applications. A line terminated in \(Z_0\)
looks infinitely long---no reflection---because the load absorbs the traveling
wave exactly as more line would.

\section{Reflection Coefficient}
At the load,
\[
\Gamma=\frac{Z_L-Z_0}{Z_L+Z_0}.
\]
If \(Z_L=Z_0\), then \(\Gamma=0\) and there is no reflection. If the load is open
or shorted, \(|\Gamma|=1\) in the lossless idealization: all the incident wave
returns.

\section{Standing Waves, SWR, and Return Loss}
Incident and reflected waves superpose into a stationary interference pattern---a
standing wave---whose voltage maxima and minima are fixed in position. Their ratio
is the standing-wave ratio,
\[
\mathrm{SWR}=\frac{V_{\max}}{V_{\min}}=\frac{1+|\Gamma|}{1-|\Gamma|}.
\]
SWR equals 1 for a perfect match and rises with mismatch. Return loss expresses
the same mismatch in decibels,
\[
RL=-20\log_{10}|\Gamma|,
\]
with higher return loss meaning a smaller reflection.

\section{Input Impedance Along the Line}
Because the reflected wave changes phase as it travels, the impedance seen
looking into a lossless line of length \(\ell\) toward a load \(Z_L\) rotates with
position:
\[
Z_{\mathrm{in}}(\ell)=Z_0\,
\frac{Z_L+\jj Z_0\tan\beta\ell}{Z_0+\jj Z_L\tan\beta\ell},
\qquad
\beta=\frac{2\pi}{\lambda}.
\]
The phase constant \(\beta\) counts radians of electrical length per unit
distance. The familiar special cases are limits of this one formula:
\begin{description}
\item[Quarter-wave (\(\beta\ell=\pi/2\), \(\tan\to\infty\))]
\(Z_{\mathrm{in}}=Z_0^2/Z_L\). The line \emph{inverts} the load impedance, which
is why \(Z_t=\sqrt{Z_0 Z_L}\) matches two real resistances.
\item[Half-wave (\(\beta\ell=\pi\), \(\tan=0\))] \(Z_{\mathrm{in}}=Z_L\): the line
repeats the load impedance.
\item[Matched (\(Z_L=Z_0\))] \(Z_{\mathrm{in}}=Z_0\) for any length: nothing to
rotate.
\end{description}
A shorted or open stub (\(Z_L=0\) or \(\infty\)) makes \(Z_{\mathrm{in}}\) a pure
reactance set by length---the basis of stub matching and line-section reactances.

\section{Quarter-Wave and Half-Wave Transformers}
From the limits above, a quarter-wave section of impedance
\(Z_t=\sqrt{Z_0Z_L}\) matches a real load \(Z_L\) to a real line \(Z_0\), and a
half-wave section repeats the load. Both assume a lossless line and hold only near
the design frequency, since \(\beta\ell\) is frequency dependent.

\section{Smith Chart Interpretation}
A Smith chart is a graphical map of normalized complex impedance and reflection
coefficient. It plots the bilinear (M\"obius) transform between normalized
impedance \(z=Z/Z_0\) and \(\Gamma=(z-1)/(z+1)\), which turns resistance and
reactance grid lines into circles and arcs inside the unit \(|\Gamma|\le1\) disk.
Crucially, the \(Z_{\mathrm{in}}(\ell)\) rotation above becomes a simple rotation
about the chart's center: moving along the line spins the point on a
constant-\(|\Gamma|\) (constant-SWR) circle. For this book the point is
conceptual---the chart is not a new law of circuits but a coordinate transform
that makes the input-impedance formula something you can read off by eye.

\begin{figure}[htbp]
\centering
\begin{tikzpicture}[scale=2.0]
\draw[thick] (0,0) circle (1);
\draw[GuideBlue,thick] (-1,0)--(1,0);
\draw[GuideLine] (0.333,0) circle (0.667);
\draw[GuideLine] (0.6,0) circle (0.4);
\draw[GuideGreen] (1,0) arc (0:180:0.5);
\draw[GuideGreen] (1,0) arc (0:-180:0.5);
\fill (0,0) circle (0.015) node[below left]{match};
\node[below] at (0,-1.12){normalized impedance plane mapped to \(\Gamma\)};
\end{tikzpicture}
\caption{Simplified Smith-chart geometry: constant-resistance circles and
reactance arcs. Moving along the line rotates a point about the center on a
constant-SWR circle.}
\end{figure}

\section{Worked Example}
A \(\SI{100}{\ohm}\) load terminates a \(\SI{50}{\ohm}\) line:
\[
\Gamma=\frac{100-50}{100+50}=\frac{1}{3},
\qquad
\mathrm{SWR}=\frac{1+1/3}{1-1/3}=2,
\qquad
RL=-20\log_{10}(1/3)\approx \SI{9.54}{\decibel}.
\]
A quarter-wave section would need
\(Z_t=\sqrt{(50)(100)}=\SI{70.7}{\ohm}\) to remove the mismatch, giving
\(Z_{\mathrm{in}}=Z_t^2/Z_L=70.7^2/100=\SI{50}{\ohm}\).
